% cv.tex — CV académico (versión en español). Build: make build (latexmk: XeLaTeX + biber)
% El contenido vive en content/*.tex (un archivo por sección).
% bib/references.bib lo sincroniza Paperpile — nunca editar a mano.
\documentclass[9pt,letterpaper]{scrartcl}

% ── Idioma ───────────────────────────────────────────────────────────────────
% es-noshorthands evita que babel active " . < > como abreviaturas, que
% chocarían con LaTeX/URLs.
\usepackage[spanish,es-noshorthands]{babel}

% ── Disposición de página ────────────────────────────────────────────────────
\usepackage[margin=0.5in,heightrounded,includefoot,footskip=0.5in]{geometry}
\usepackage{parskip}

% ── Tipografías (XeLaTeX) ────────────────────────────────────────────────────
\usepackage{fontspec}
\usepackage{fontawesome5}

% ── Tablas ───────────────────────────────────────────────────────────────────
\usepackage{longtable}
\usepackage{array}

% ── Colores y enlaces ────────────────────────────────────────────────────────
\usepackage{xcolor}
\definecolor{linkcolour}{rgb}{0,0.2,0.6}
\definecolor{darkgray}{gray}{0.4}
\usepackage{hyperref}
\hypersetup{colorlinks=true, urlcolor=linkcolour, linkcolor=linkcolour}

% ── Formato de secciones ─────────────────────────────────────────────────────
% \MakeUppercase (no el primitivo \uppercase) para que las vocales acentuadas
% de los títulos se conviertan bien: "Formación" → "FORMACIÓN".
\usepackage{titlesec}
\titleformat{\section}
  {\normalfont\normalsize\bfseries\raggedright\MakeUppercase}
  {\thesection}{1em}{}
\titleformat{\subsection}
  {\normalfont\normalsize\bfseries\raggedright}
  {\thesubsection}{1em}{}
\titleformat{\subsubsection}
  {\normalfont\normalsize\bfseries\raggedright}
  {\thesubsubsection}{1em}{}
\setcounter{secnumdepth}{0}   % CV académico: sin numeración de secciones

% ── Bibliografía (reemplaza read_bib/format_*_vancouver de preprocess.R) ─────
\usepackage[backend=biber,style=vancouver,maxnames=6,minnames=6,
            defernumbers=true,sorting=none]{biblatex}
\addbibresource{bib/references.bib}

% Correcciones de salida Vancouver/ICMJE (ver revisión de la migración):
% - quitar language ("en.") y urldate ("[Accessed on: ...]") de las entradas
% - coma (no "y") antes del último autor
% - bibliografía en bandera derecha: evita que el justificado estire el DOI/URL
\AtEveryBibitem{\clearlist{language}\clearfield{langid}\clearfield{urldate}\clearfield{urlday}\clearfield{urlmonth}\clearfield{urlyear}\clearfield{urldateera}}
% spanish.lbx cambia la última coma por " y "/" e " mediante el contador
% smartand y su propio delimitador, aplicado al cargar el idioma. Se apaga
% smartand y se fija el delimitador final a coma dentro de \AtBeginDocument
% (después de que el idioma se cargue) para forzar el estilo Vancouver.
\setcounter{smartand}{0}
\AtBeginDocument{\DeclareDelimFormat{finalnamedelim}{\addcomma\space}}
\appto{\bibsetup}{\raggedright}

% Categorías (reemplaza data/pub_categories.yaml + el registro section_labels).
% \DeclareBibliographyCategory es solo de preámbulo; la asignación de claves
% ocurre en content/publications.tex vía \cvpub.
\DeclareBibliographyCategory{peerreviewed}
\DeclareBibliographyCategory{conference}
\DeclareBibliographyCategory{thesis}
\DeclareBibliographyCategory{whitepaper}
\DeclareBibliographyCategory{preprints}
\DeclareBibliographyCategory{bookchapter}
\newcommand{\cvpub}[2]{\addtocategory{#1}{#2}\nocite{#2}}

% Campos que Paperpile nunca trae (reemplaza bib_patches/patch_bib).
\DeclareSourcemap{%
  \maps[datatype=bibtex]{%
    \map{%
      \step[fieldsource=entrykey, match={Vidyasagaran2024-wx}, final]
      \step[fieldset=addendum,
            fieldvalue={Entre los autores: Del Castillo Fernandez, D.}]%
    }%
  }%
}

% Resaltar el nombre propio en negrita (reemplaza filters/bold-author.lua).
% Misma regla: apellido que contiene "Del" y "Castillo", con nombre de pila
% que empieza por "D" (cubre "Del Castillo", "Del Castillo Fernández",
% "Del Castillo-Fernández").
\usepackage{xstring}
\newif\ifownname
\newcommand*{\checkownname}{%
  \global\ownnamefalse
  \ifdefvoid{\namepartgiven}{}{%
    \ifdefvoid{\namepartfamily}{}{%
      \begingroup\expandarg
        % \namepartfamily es 'Del\bibnamedelima Castillo' — hay una macro
        % delimitadora entre los tokens, así que un prefijo 'Del Castillo'
        % nunca coincide. Las pruebas de subcadena la saltan.
        \IfSubStr{\namepartfamily}{Castillo}
          {\IfSubStr{\namepartfamily}{Del}
            {\IfBeginWith{\namepartgiven}{D}{\global\ownnametrue}{}}
            {}}%
          {}%
      \endgroup
    }%
  }%
}
\renewcommand*{\mkbibnamefamily}[1]{\checkownname\ifownname\textbf{#1}\else#1\fi}
\renewcommand*{\mkbibnamegiven}[1]{\checkownname\ifownname\textbf{#1}\else#1\fi}

\newcommand{\myorcid}{0000-0002-8609-0312}

\begin{document}
\noindent
\begin{minipage}[c]{0.50\textwidth}
{\huge\underline{\textbf{Darwin Del Castillo, MD, MPH}}}
\end{minipage}%
\hfill
\begin{minipage}[c]{0.50\textwidth}
\raggedleft\large
\faOrcid\hspace{0.4em}\href{https://orcid.org/0000-0002-8609-0312}{0000-0002-8609-0312} \\[2pt]
\faGithub\hspace{0.4em}\href{https://github.com/ddelcastillof}{ddelcastillof}\\[2pt]
\faGlobe\hspace{0.4em}\href{https://www.gla.ac.uk/schools/healthwellbeing/staff/darwindelcastillofernandez/}{UofG Profile}
\end{minipage}
\vspace{4pt}

\section{Education}
\vspace{-1.5em}
\textcolor{darkgray}{\rule{\textwidth}{0.5pt}}
\begin{longtable}{p{14cm}>{\raggedleft\arraybackslash}p{4cm}}
\textbf{Master of Public Health (MPH)} & \textbf{Sep 2023 -- Jun 2025} \\
\textsc{University of Washington} & \textsc{Seattle -- USA} \\
& \\
\textbf{Doctor of Medicine (MD)} & \textbf{Mar 2014 -- Aug 2021} \\
\textsc{Universidad Nacional Mayor de San Marcos} & \textsc{Lima -- Peru} \\
& \\
\end{longtable}

\section{Additional Education}
\vspace{-1.5em}
\textcolor{darkgray}{\rule{\textwidth}{0.5pt}}
\begin{longtable}{p{14cm}>{\raggedleft\arraybackslash}p{4cm}}
\textbf{Professional Certificate in Clinical Research} & \textbf{Sep 2021 -- Apr 2022} \\
\textsc{Universidad Peruana Cayetano Heredia} & \textsc{Lima -- Peru} \\
& \\
\end{longtable}

\section{Experiencia Profesional}
\vspace{-1.5em}
\textcolor{darkgray}{\rule{\textwidth}{0.5pt}}
\begin{longtable}{p{14cm}>{\raggedleft\arraybackslash}p{4cm}}
\textbf{Asistente de Investigación} & \textbf{Dic 2025 -- Actualidad} \\
\textsc{School of Health and Wellbeing -- University of Glasgow} & \textsc{Glasgow -- UK} \\
\multicolumn{2}{p{18cm}}{%
\begin{itemize}
  \item Estimé los efectos causales del desempleo sobre desenlaces de salud física y mental empleando métodos-g con datos administrativos.
  \item Realicé análisis de mediación causal del ingreso relativo entre el desempleo y la mortalidad por todas las causas usando datos administrativos enlazados.
  \item Modelé el efecto de implementar políticas económicas sobre desenlaces de salud mediante modelamiento por microsimulación.
\end{itemize}
} \\
\\
\textbf{Asistente de Investigación de Posgrado} & \textbf{Jul 2024 -- Jun 2025} \\
\textsc{Institute for Health Metrics and Evaluation} & \textsc{Seattle -- USA} \\
\multicolumn{2}{p{18cm}}{%
\begin{itemize}
  \item Desarrollé hiperparámetros para un modelo de meta-regresión jerárquico bayesiano con el fin de estimar prevalencia, incidencia y ocurrencia de desenlaces clínicos de infecciones del tracto respiratorio inferior en 80 países.
  \item Ejecuté cadenas MCMC utilizando priors basados en datos y propagación de la incertidumbre para estimar el traslape probabilístico entre grupos de inclusión social en 22 países de altos ingresos.
  \item Estimé fracciones atribuibles poblacionales y razones de mortalidad estandarizadas a partir del traslape entre grupos de inclusión social.
\end{itemize}
} \\
\\
\textbf{Asistente de Investigación de Posgrado} & \textbf{Jun 2024 -- Mar 2025} \\
\textsc{Department of Global Health -- University of Washington} & \textsc{Seattle -- USA} \\
\multicolumn{2}{p{18cm}}{%
\begin{itemize}
  \item Diseñé un protocolo de estudio para evaluar la efectividad de estrategias de implementación para la atención de quemaduras pediátricas en establecimientos de salud peruanos.
  \item Coordiné un estudio de métodos mixtos para explorar barreras y facilitadores en más de 10 hospitales peruanos de bajos recursos.
  \item Co-desarrollé estrategias de comunicación con 2 organizaciones de actores clave para difundir hallazgos complejos entre tomadores de decisión en salud.
\end{itemize}
} \\
\\
\textbf{Asistente de Investigación de Posgrado} & \textbf{Sep 2023 -- Mar 2024} \\
\textsc{Health Population Research Center -- University of Washington} & \textsc{Seattle -- USA} \\
\multicolumn{2}{p{18cm}}{%
\begin{itemize}
  \item Mantuve y almacené datos sociodemográficos y de utilización como parte de una encuesta de seguimiento posterior a un ensayo clínico en escolares, utilizando plataformas electrónicas de captura de datos.
  \item Traduje materiales de sometimiento al comité de ética (IRB) del español al inglés para el Alzheimer's Disease Research Center.
  \item Analicé datos de participación y alcance de ensayos clínicos para generar herramientas de visualización y tablas resumen, permitiendo que los gestores de programa ajustaran las estrategias de reclutamiento.
\end{itemize}
} \\
\\
\newpage
\textbf{Consultor de Investigación} & \textbf{Ene 2023 -- Ene 2024} \\
\textsc{Health Technology Assessment and Research Institute} & \textsc{Lima -- Perú} \\
\multicolumn{2}{p{18cm}}{%
\begin{itemize}
  \item Integré datos de vigilancia y datos clínicos de hospitales peruanos en modelos estadísticos de carga de enfermedad para apoyar los procesos nacionales de evaluación de tecnologías sanitarias y priorización.
  \item Realicé análisis de impacto presupuestario y modelos de costo-efectividad que priorizaron intervenciones oncológicas por montos superiores a USD 1 millón.
  \item Sinteticé los hallazgos en informes ejecutivos dirigidos a tomadores de decisión y socios, orientando la política de reembolso del Seguro Social de Salud del Perú.
\end{itemize}
} \\
\\
\textbf{Asistente de Investigación} & \textbf{Mar 2019 -- Ago 2023} \\
\textsc{CRONICAS Center of Excellence in Chronic Diseases} & \textsc{Lima -- Perú} \\
\multicolumn{2}{p{18cm}}{%
\begin{itemize}
  \item Coordiné 2 talleres por sede con agentes comunitarios de salud en 3 regiones peruanas para co-diseñar una aplicación móvil para la detección temprana del deterioro cognitivo.
  \item Contribuí a propuestas de financiamiento en epidemiología de enfermedades no transmisibles en poblaciones pediátricas y adultas mediante análisis preliminares de datos clínicos y demográficos, apoyando el desarrollo de proyectos de investigación.
  \item Fui coautor de 5 manuscritos científicos, incluyendo análisis de datos longitudinales, revisiones sistemáticas y modelos de aprendizaje automático.
\end{itemize}
} \\
\\
\end{longtable}

\section{Skills}
\vspace{-1.5em}
\textcolor{darkgray}{\rule{\textwidth}{0.5pt}}
\textbf{Programming languages:} R, Python, Stata, SQL (MySQL) \\
\textbf{Cloud Computing:} Slurm HPC, Containerization (Docker/Apptainer), Linux environments (bash) \\
\textbf{Tools:} VS Code, Quarto, \LaTeX, Git version control (GitHub, Bitbucket) \\
\textbf{Languages:} English (fluent), Spanish (native) \\
\textbf{Core Competencies:} Data analysis, teamwork, problem solving, science communication, mentorship \\

\section{Distinciones Académicas y Premios}
\vspace{-1.5em}
\textcolor{darkgray}{\rule{\textwidth}{0.5pt}}
\begin{longtable}{>{\raggedright\arraybackslash}p{16.5cm} p{1.5cm}}
Beca para el Taller Metodológico del Estudio Young Lives, \textsc{Grupo de Análisis para el Desarrollo (GRADE)} -- Perú & \textbf{Jun 2018} \\
& \\
Graduación de la Escuela de Medicina con Honores \textit{(Cum Laude)} (Promedio: 17/20), \textsc{Facultad de Medicina, Universidad Nacional Mayor de San Marcos} -- Perú & \textbf{Ago 2021} \\
& \\
Beca Fogarty para el Diplomado en Investigación Clínica, \textsc{Escuela de Salud Pública, Universidad Peruana Cayetano Heredia} -- Perú & \textbf{Sep 2021} \\
& \\
Beca del Departamento de Salud Global por Excelencia, Equidad e Impacto en Salud Global, \textsc{Department of Global Health, University of Washington} -- USA & \textbf{Feb 2024} \\
& \\
\end{longtable}

\section{Teaching Experience}
\vspace{-1.5em}
\textcolor{darkgray}{\rule{\textwidth}{0.5pt}}

\subsection{Teaching Contributions}
\begin{longtable}{@{}p{2.5cm}>{\raggedright\arraybackslash}p{16cm}@{}}
\textbf{2025 -- Current} & \textbf{Instructor} -- Introduction to Epidemiology \\*
 & \textsc{Universidad Científica del Sur, Lima -- Peru} \\*
& \\
\textbf{2024} & \textbf{Teaching Assistant} -- Introduction to Molecular and Cellular Biology \\*
 & \textsc{University of Washington, Seattle -- USA} \\*
\end{longtable}
%%% include supervised thesis dissertations
\section{Research Contributions \hspace{1cm} ORCID: \myorcid}
\vspace{-1.5em}
\textcolor{darkgray}{\rule{\textwidth}{0.5pt}}

% New-paper workflow: Paperpile bot pushes bib/references.bib; add ONE
% \cvpub line below (order = display order); make build.
% Unknown cite key => biber warning in build/cv.blg and make check fails.
\cvpub{peerreviewed}{Del-Castillo-Fernandez2021-kl}
\cvpub{peerreviewed}{Lapidaire2023-pr}
\cvpub{peerreviewed}{Vidyasagaran2024-wx}
\cvpub{peerreviewed}{Bazo-Alvarez2026-yz}
\cvpub{peerreviewed}{Pinzon-Cortes2026-mv}
\cvpub{peerreviewed}{Pinzon-Cortes2026-vu}

% Conference abstract. Paperpile records it as an @article, so the only thing
% marking it as an abstract is the addendum patched in cv.tex's
% \DeclareSourcemap — see the comment on the Del-Castillo2026-uy map.
\cvpub{conference}{Del-Castillo2026-uy}

\printbibliography[category=peerreviewed,title={Journal Manuscripts},
                   heading=subbibliography,resetnumbers=true]

\printbibliography[category=conference,title={Conference Abstracts},
                   heading=subbibliography,resetnumbers=true]

% Categories below are empty today. When the first key lands, add its
% \cvpub line above and uncomment the matching \printbibliography:
% \printbibliography[category=thesis,title={Thesis and Dissertations},heading=subbibliography,resetnumbers=true]
% \printbibliography[category=whitepaper,title={White Papers},heading=subbibliography,resetnumbers=true]
% \printbibliography[category=preprints,title={Preprints},heading=subbibliography,resetnumbers=true]
% \printbibliography[category=bookchapter,title={Book Chapters},heading=subbibliography,resetnumbers=true]

\subsection{Peer-Review Contributions}
\begin{tabular}{lll}
Peer Reviewer & SSM - Population Health & 2025 \\
\end{tabular}

\section{Licensure and Certification}
\vspace{-1.5em}
\textcolor{darkgray}{\rule{\textwidth}{0.5pt}}
\begin{longtable}{@{}p{2.5cm}>{\raggedright\arraybackslash}p{16cm}@{}}
\textbf{2024} & National Registry of Investigators \textsc{(RENACYT-CONCYTEC)}, Investigator Level VII (P0160977) \\
\textbf{2023} & CITI Program, Good Clinical Practice and ICH, Credential ID 65079238 \\
\textbf{2023} & CITI Program, Human Subjects Learners, Credential ID 65079240 \\
\textbf{2023} & CITI Program, Biomedical Responsible Conduct of Research, Credential ID 65079239 \\
\textbf{2021} & Peruvian Medical Board, License No. 94534 \\
\end{longtable}

\section{Afiliaciones a Organizaciones Profesionales}
\vspace{-1.5em}
\textcolor{darkgray}{\rule{\textwidth}{0.5pt}}
\begin{itemize}
 \item John Snow Society (2026 -- Actualidad)
 \item Society for Social Medicine \& Population Health (2026 -- Actualidad)
 \item International Epidemiological Association (2020 -- Actualidad)
 \item Colegio Médico del Perú (2021 -- Actualidad)
 \item Consortium of Universities for Global Health (2023 -- Actualidad)
 \item America's Network for Chronic Disease Surveillance (2025 -- Actualidad)
 \item Global Burden of Disease Study Collaborator (2025 -- Actualidad)
\end{itemize}

\end{document}
